\section{Which gradient reaches an interior core?}
\label{sec:interior}
The root calculation extends to every internal core of an aligned multilinear
student under squared loss. An interior core receives information from both
the child subtrees and the part of the network that uses its output.
Conditioning on its grammar symbol separates those two contributions.

At position $i$ put $z_i=h_{iL}\otimes h_{iR}$, $h_i=W_i z_i$,
and write the student output as $f=A_i h_i$. Here $A_i$ maps changes
at this occurrence of the bond to the logits, holding the other subtree
activations fixed. It depends only on observations outside the subtree at $i$
when parameters are fixed. Denote the grammar symbol by $S_i=s$,
with probability $p_i(s)>0$, and define
\[
 d_i(s)=\mathbb E[z_i\mid s],\quad F_i(s)=\mathbb E[z_i z_i^\top\mid s],
 \quad m_i(s)=\mathbb E[A_i^\top e_Y\mid s],\quad
 N_i(s)=\mathbb E[A_i^\top A_i\mid s].
\]
Then the negative gradient from this occurrence is exactly
\begin{equation}
 -\nabla_{W_i}\mathcal L_2
 =\sum_s p_i(s)\left[m_i(s)d_i(s)^\top-N_i(s)W_iF_i(s)\right].
 \label{eq:interior}
\end{equation}
Indeed the sample gradient is
$A_i^\top(e_Y-A_i W_i z_i)z_i^\top$ with the descent sign.
Conditioned on $s$, the outside variables $(A_i,Y)$ are independent of
the draws that expand this subtree. Taking the expectation of each product
yields~\eqref{eq:interior}. If several occurrences share a table, sum all
their partial derivatives. Freezing other occurrences while evaluating
each partial derivative is precisely the chain rule for tying.

The two inward messages $d_i,F_i$ are already available in
Section~\ref{sec:categorical-root}. The two outward messages also have an
exact recursion. At the root use $A=I$, $m(y)=e_y$, $N(y)=I$.
For a left descent from parent $i$, let $B_i(v)$ denote the matrix
mapping $u$ to $W_i(u\otimes v)$. For fixed sibling activation $h_R$,
$A_L=A_i B_i(h_R)$. The probability that the parent is $a$ and the
sibling symbol is $c$, conditional on the left symbol $b$, is
\[
 \omega(a,c\mid b)=\frac{p_i(a)G_{\ell(i)}(a,b,c)}
 {\sum_{a',c'}p_i(a')G_{\ell(i)}(a',b,c')}.
\]
Unused symbols can be skipped. Writing $\mu_R(c)=\mathbb E[h_R\mid c]$,
independence of the outer context and the parent production conditional
on $a$, and independence of the sibling expansion given $c$, give
\begin{align*}
 m_L(b)&=\sum_{a,c}\omega(a,c\mid b)\,
        B_i(\mu_R(c))^\top m_i(a),\\
 N_L(b)&=\sum_{a,c}\omega(a,c\mid b)\,
 \mathbb E[B_i(h_R)^\top N_i(a) B_i(h_R)\mid c].
\end{align*}
The last integrand is quadratic in $h_R$, so its conditional second
moment suffices. Exchanging child roles gives the corresponding right
messages. At the root, \eqref{eq:interior} reduces to $D/n-CF$ with the
notation of \eqref{eq:full-square}. This argument needs neither independent
child symbols marginally nor a diagonal student. It gives all population
gradients without first selecting a set of teacher modes. A particular
update can still cancel after summing over symbols or tied positions.
For softmax loss the predicted-probability term depends on the entire
logit vector, so this squared-loss factorization no longer follows.
